%------------------------------------------------------------------------------%
\begin{question}
  Dados
  \[
    \vec{w} = \vetortri{2}{-1}{3}
    \qquad
    \vec{v} = \vetortri{1}{2}{2}
  \]
  Determine a decomposição
  \[
    \vec{w} = \vec{w}_{\parallel} + \vec{w}_{\perp}
  \]

\end{question}

%------------------------------------------------------------------------------%
\begin{resolution}{Solução}
  Do exercício anterior temos
  \pause
  \[
    \vec{w}_{\parallel}
    \pause
    \;=\; \projuv{w}{v}
    \pause
    \;=\; \Vetortri{\frac{2}{3}}{\frac{4}{3}}{\frac{4}{3}}
  \]
\end{resolution}

%------------------------------------------------------------------------------%
\begin{resolution}{Solução}
  \begin{align*}
    \onslide<+->{\vec{w}_{\perp}}
    \onslide<+->{\;=\; \vec{w} - \vec{w}_{\parallel}}
    \onslide<+->{& =   \vetortri{2}{-1}{3} - \Vetortri{\frac{2}{3}}{\frac{4}{3}}{\frac{4}{3}}}
    \onslide<+->{\;=\; \Vetortri{2-\frac{2}{3}}{-1-\frac{4}{3}}{3-\frac{4}{3}} \\}
    \onslide<+->{& =   \Vetortri{\frac{6-2}{3}}{\frac{-3-4}{3}}{\frac{9-4}{3}}}
    \onslide<+->{\;=\; \Vetortri{\frac{4}{3}}{-\frac{7}{3}}{\frac{5}{3}}}
  \end{align*}
\end{resolution}

%------------------------------------------------------------------------------%
\endinput